\section{A relative prime number theorem for the window moduli}
\label{sec:window-relative-pnt}

The denominator in Section~\ref{sec:window-exceptions} counts primes
in a single reduced residue class. An exceptional real zero can make
this count small. We prove the estimate with an error proportional to
the corrected main term, and then give the lower bound on the
denominator for the moduli used in the two selections of primes.

For \(z\ge2\) set
\(\ell(z)=\int_2^z dt/\log t\), and let \(\pi(X;M,b)\) count primes
\(p\le X\) congruent to \(b\bmod M\). Call \((\chi,\beta)\) eligible
at \(M\) if \(\chi\) is a nonprincipal real Dirichlet character modulo
\(M\), \(L(\beta,\chi)=0\), and
\[
       1-\frac{1}{8\log M}<\beta<1.
\]
There are absolute positive constants \(c,A,C\) with the following
property. For every positive multiple \(M\) of \(8\), select one
eligible pair if one exists, and select the empty option otherwise.
The selection can be made depending only on \(M\), in such a way that
for all integers \(b\) coprime to \(M\) and all real \(X\ge2\) with
\[
             \log X\ge C\log M\,\log(\log M+2)
\]
one has
\begin{align}
 B(X)&=\frac{\ell(X)-\chi(b)\ell(X^\beta)}{\varphi(M)},&
 E_c(X,M)&=\exp\!\left(-\frac{c\log X}{\log M}\right)
                   +\exp\!\left(-\sqrt{c\log X}\right),
 \nonumber\\
 |\pi(X;M,b)-B(X)|&\le A E_c(X,M) B(X).
 \label{eq:window-relative-pnt}
\end{align}
The character term is omitted for the empty option. The choice of
\(C\) below ensures \(X^\beta\ge2\) everywhere the estimate is used.

\paragraph{The weighted count.}
We use the unconditional progression theorem of J.~Thorner and
A.~Zaman. Their exceptional pair has \(\beta\ge1-1/(50\log M)\).
McCurley's zero-free region gives uniqueness in this interval,
simplicity, and a nonprincipal real character. With the exceptional
pair in their sense put \(\sigma=\chi(b)\) and
\(F(t)=t-\sigma t^\beta/\beta\). In its absence put \(F(t)=t\)
and \(\sigma=0\). Define
\(\vartheta(t;M,b)=\sum_{p\le t,\ p\equiv b\bmod M}\log p\).
When \(M\ge2\), \((b,M)=1\), \(t\ge3\), and
\(F(t)/\varphi(M)\ge t^{18/19}\), their Corollary~1.5 gives
\[
 \left|\vartheta(t;M,b)-\frac{F(t)}{\varphi(M)}\right|
 \ll \frac{F(t)}{\varphi(M)}
 \left\{\exp\!\left(-\frac{a\log t}{\log M}\right)
       +\exp\!\left(-a\frac{(\log t)^{3/5}}{(\log\log t)^{1/5}}\right)
 \right\}
\]
for an absolute \(a>0\). We use the weaker consequence
\(1-\beta\gg M^{-1}\) of the exceptional-zero estimate in their
equation~(1.2). Both results concern all characters at modulus
\(M\), which permits imprimitive characters. Equivalently, the
extra Euler factors of an induced character are nonzero at real
\(\beta>0\), since \(|\chi^*(p)p^{-\beta}|<1\).
See Thorner and Zaman, \emph{Refinements to the prime number theorem
for arithmetic progressions}, Math.~Z. 306 (2024), article~54,
\href{https://doi.org/10.1007/s00209-023-03414-3}{journal publication},
\href{https://arxiv.org/abs/2108.10878}{authors' text}.

Hereafter increase the absolute \(C\) so that
\(s=\sqrt X\ge M^{100}\) and \(s\) exceeds any required absolute
threshold. This is possible since \(M\ge8\).
For the sign \(\sigma=1\) write \(\delta=1-\beta\) and
\(w(t)=\min(1,\delta\log t)\). For \(\sigma=0,-1\) put \(w(t)=1\).
Uniformly on \(t\ge s\) we have \(F(t)\asymp t w(t)\).
For \(\sigma=1\), indeed, \(\beta>1/2\) yields
\(0\le-\log\beta\le2\delta\). Thus for \(\log t\ge4\),
\[
   1-e^{-\delta\log t/2}
      \le \frac{F(t)}{t}
      =1-\frac{e^{-\delta\log t}}{\beta}
      \le 1-e^{-\delta\log t}.
\]
Use \(1-e^{-v}\asymp\min(1,v)\) for \(v\ge0\).
For \(\sigma=-1\) the claimed comparison follows from
\(1\le F(t)/t\le3\). In all three cases \(w(t)\gg M^{-1}\).
Since \(\varphi(M)\le M\),
\[
           F(t)/\varphi(M)\gg t M^{-2}\ge t^{49/50}.
\]
This exceeds \(t^{18/19}\) once \(s\) exceeds an absolute constant.
The published hypothesis therefore holds at every \(t\in[s,X]\),
for the same class and the same exceptional option.

\paragraph{Removing the logarithmic weight.}
First form \(B\) with the authors' exceptional option.
For \(\sigma=1\),
\[
  \varphi(M) B(X)=\int_{X^\beta}^{X}\frac{dt}{\log t}
     \ge \frac{X(1-e^{-\delta\log X})}{\log X}.
\]
For \(\sigma=0,-1\), integration over \([X/2,X]\)
gives the corresponding lower bound \(X/(2\log X)\).
Together these inequalities yield
\begin{equation}
 B(X)\gg \frac{X w(X)}{\varphi(M)\log X}
          \gg\frac{X}{M^2\log X}.
 \label{eq:window-ap-main-lower}
\end{equation}
Write \(H(t)=\vartheta(t;M,b)-F(t)/\varphi(M)\).
Uniformly for \(s\le t\le X\), the weighted theorem gives
\[
 |H(t)|\ll \frac{t w(X)}{\varphi(M)} E_c(X,M)
\]
after choosing \(c>0\) sufficiently small. Indeed
\(\log t\ge\tfrac12\log X\), \(w(t)\le w(X)\), and
\(u^{3/5}/(\log u)^{1/5}\gg\sqrt u\) for large \(u\).
On this interval \(\varphi(M)B'(t)=F'(t)/\log t\), including the
factor \(1/\beta\) in \(F\). Partial summation with lower endpoint
\(s\) gives the exact identity
\[
 \pi(X;M,b)-B(X)=\pi(s;M,b)-B(s)
   +\frac{H(X)}{\log X}-\frac{H(s)}{\log s}
   +\int_s^X\frac{H(t)}{t(\log t)^2}\,dt.
\]
It follows either by integration by parts for the prime-mass step
function or by applying ordinary integration by parts separately to
\(F'(t)/\log t\). In particular the normalization of \(\ell\) at
\(2\) is retained at the lower endpoint. The two boundary terms
in \(H\) and the integral together are
\(O(X w(X) E_c(X,M)/(\varphi(M)\log X))\), since
\(\log s=\tfrac12\log X\) and
\(\int_s^Xdt/(\log t)^2\le4X/(\log X)^2\).
Also \(|\pi(s;M,b)-B(s)|\ll s\), by the elementary upper bounds
\(\ell(z)\ll z\) and \(\pi(s;M,b)\le s\).
As \(M\le X^{1/200}\), \eqref{eq:window-ap-main-lower} gives
\(s/B(X)\ll X^{-49/100}\log X\ll\exp(-\sqrt{c\log X})\).
This proves \eqref{eq:window-relative-pnt} for the authors' option.

\paragraph{The release's wider exceptional interval.}
If the authors' pair exists, it is eligible, so choose it. Otherwise,
if there is an eligible pair, fix any such \((\chi,\beta)\) once for
this \(M\). Its distance from \(1\) satisfies
\(1-\beta>1/(50\log M)\). Uniformly for \(\beta>1/2\) and
\(X^\beta\ge4\),
\[
 0\le\frac{\ell(X^\beta)}{\ell(X)}
     \ll X^{\beta-1}
     \le \exp\!\left(-\frac{\log X}{50\log M}\right).
\]
Here \(\ell(z)\asymp z/\log z\) for \(z\ge4\), obtained by splitting
the upper estimate at \(\sqrt z\) for large \(z\) and integrating
over \([z/2,z]\) for the lower estimate. Increase \(C\) so the last
bound, with its implied constant, is at most \(1/2\).
The newly selected main term is then between one half and twice
\(\ell(X)/\varphi(M)\), and their difference is bounded by a constant
times the new main term times \(E_c(X,M)\) on taking \(c\le1/50\).
If no eligible pair exists the empty option already works.
This proves the asserted choice for all \(M,X,b\) and also preserves
the lower bound \(B(X)\gg X/(M^2\log X)\). In particular the proof
requires uniqueness only for zeros in the narrower interval.

\paragraph{The actual prime interval.}
For every \(\epsilon>0\) there exists \(Y_0(\epsilon)\) such that
for every \(Y\ge Y_0(\epsilon)\), every \(m>0\) divisible by \(8\),
every \(a_0\) coprime to \(m\), and every \(V>0\) satisfying
\[
  \tfrac14(\log Y)^2\log m\le\log V,\qquad
  Y^{\epsilon/6}\le \log V\le Y^{\epsilon/5}(\log Y)^2,
\]
the set \(\mathcal P=\{p\text{ prime}:\sqrt V<p\le V,\
p\equiv a_0\bmod m\}\) satisfies
\begin{equation}
                       |\mathcal P|\ge V^{9/10}.
 \label{eq:window-ap-lower-count}
\end{equation}
Put \(L=\log Y\), \(u=\log V\).
Then \(\log m\le4u/L^2\) and \(\log m\le4Y^{\epsilon/5}\), whence
\(\log(\log m+2)=O_\epsilon(L)\).
Thus \(C\log m\log(\log m+2)\le u\) for sufficiently large \(Y\).
Both exponentials in \(E_c(V,m)\) tend to zero uniformly, because
\(u/\log m\ge L^2/4\) and \(u\ge Y^{\epsilon/6}\).
Consequently \(\pi(V;m,a_0)\ge B(V)/2\). For large \(Y\) one has
\(m\le V^{1/100}\). The preceding main-term lower bound gives
\(\pi(V;m,a_0)\gg V^{49/50}/\log V\). Subtracting at most
\(\sqrt V\) primes leaves at least \(V^{9/10}\) for sufficiently
large \(V\), uniformly in \(a_0,m\). The lower hypothesis on \(u\)
ensures that this last threshold follows from \(Y\ge Y_0(\epsilon)\).
The formal statement also quantifies over nonpositive real \(V\).
Under its real-power convention \(V^{9/10}=|V|^{9/10}\cos(9\pi/10)\)
for \(V<0\), and \(0^{9/10}=0\). These cases satisfy the inequality
by nonnegativity of the cardinality.

There is also uniformity over the enlarged moduli occurring in
the character sums. Given fixed \(D,B_0>0\), take \(M=ms\)
with any positive integer \(s\le Y^D\). Then
\(\log M\le4u/L^2+DL\le8u/L^2\) once \(Y\) is large, since
\(Y^{\epsilon/6}/L^3\) tends to infinity. The upper bound
\(\log M\le4Y^{\epsilon/5}+DL\) still gives
\(\log(\log M+2)=O_{\epsilon,D}(L)\).
Thus \eqref{eq:window-relative-pnt} applies with
\(E_c(V,M)=o(Y^{-B_0})\), uniformly over \(s\) and every reduced
class modulo \(M\). In particular one can use \(D=800K\) and
\(B_0=K+1\) for fixed \(K\ge1\) in Section~\ref{sec:window-exceptions}.

\paragraph{Exceptional pairs at two different moduli.}
For the use of Landau--Page uniqueness in
Section~\ref{sec:window-exceptions}, let \(H\ge3\), take integers
\(1\le M_1,M_2\le H\), and suppose nonprincipal real characters
\(\chi_i\bmod M_i\) have real zeros
\[
                  1-\frac{1}{100\log H}<\beta_i<1 .
\]
Put \(q=M_1M_2\). The induced character \(\widetilde\chi_i\bmod q\)
is nonprincipal, since the reduction map on unit groups is
surjective. This surjectivity follows by lifting a unit modulo
\(M_i\) to an integer congruent to \(1\) at each prime divisor of
\(q\) absent from \(M_i\), using the Chinese remainder theorem.
The identity
\[
 L(s,\widetilde\chi_i)=L(s,\chi_i)
   \prod_{\substack{p\mid q\\p\nmid M_i}}(1-\chi_i(p)p^{-s})
\]
preserves the zero at \(\beta_i\).
Since \(\log q\le2\log H\), both zeros lie in
\([1-1/(50\log q),1)\). Fixed-modulus uniqueness at \(q\) now yields
\(\widetilde\chi_1=\widetilde\chi_2\) and \(\beta_1=\beta_2\).
In particular \(\chi_1(n)=\chi_2(n)\) for all integers \(n\)
coprime to \(M_1M_2\). This supplies precisely the cross-modulus
conclusion of \texttt{MSqPage.page\_across} and verifies that the
same primitive character accounts for all such zeros. The constant
\(1/100\) is independent of the two moduli and the window. The
fixed-modulus theorem is due to K.~S.~McCurley,
\href{https://doi.org/10.1016/0022-314X(84)90089-1}{\emph{Explicit
zero-free regions for Dirichlet \(L\)-functions}}, J.~Number Theory
19 (1984), 7--32, in the form stated by Thorner and Zaman.

\paragraph{Source correspondence.}
In the tagged release the definitions of \(\ell,\pi,B\) and
\srcid{PNTAPRelW} are in \srcid{MeanSquareHypWIff}. The declaration
\texttt{pnt\_progressions\_relative\_weak\_holds} in
\srcid{PNTFinal} proves the estimate internally. Its
\srcid{PNTGlue} separates the ranges
\(\log M\le c_A\sqrt{\log t}\) and \(\log M>c_A\sqrt{\log t}\).
The latter uses the smoothed explicit formula, a differenced
explicit formula, the zero-density estimate, and the
Deuring--Heilbronn estimate. \srcid{PrimeCountOfPsi} converts the
combined estimate to the unweighted prime count. The proof above
uses the published weighted theorem to establish the same
\srcid{PNTAPRelW} directly. The window count
\eqref{eq:window-ap-lower-count} is \texttt{MSqCount.count\_a}.
These modules are in
\srcid{Formalism.GL3MeanSquares.GenusMeanAsymptotics.JacobiCharacterSums}.
The common-window application is
\srcid{Development.AutomorphicArithmetic.MeanSquareBounds.twistK}.
